\chapter{Production-Readiness and Governance Checklists}
\index{production readiness}\index{checklists}\index{governance!checklist}%
This appendix distills the chapters' production notes, acceptance criteria, and operational
guidance into pre-launch checklists, one per domain. They are written as review artifacts: walk
the relevant list with the team before a workload is declared production, and tick items with
evidence, not intention. Chapter references point to the source material.

\providecommand{\chk}{\item[$\square$]}

\section{Storage and Data Lake (S3, Chapter 1)}
\index{checklist!storage}%
\begin{itemize}
\chk Versioning enabled on every bucket that holds data you cannot regenerate.
\chk Lifecycle rules transition and expire per the data's retention class; the cost model was
validated with the actual object-size distribution (small objects can cost \emph{more} after a
Glacier transition).
\chk Object Lock mode chosen deliberately: governance mode only when an administrative bypass
process is documented; compliance mode only when retention is certain.
\chk Default encryption is SSE-KMS with a key policy that names the allowed principals; public
access blocked at the account level.
\chk Replication (CRR/SRR) configured where the RPO demands it, with a replication-time alarm.
\chk Zone naming convention (\texttt{raw/}, \texttt{validated/}, \texttt{curated/}) enforced by
bucket or prefix policy, and every prefix has an owner tag.
\chk Event notifications (SQS/SNS/EventBridge) tested end-to-end with a poison-message path.
\chk Deletion runbook exists and has been rehearsed; nothing relies on ``we will remember.''
\end{itemize}

\section{Warehousing and Databases (RDS/Aurora, DynamoDB, Redshift; Chapters 2, 3, 5, 6)}
\index{checklist!warehousing}%
\begin{itemize}
\chk Multi-AZ for the primary; failover rehearsed, not just configured (Chapter~2 lab pattern).
\chk PITR window matches the recovery objective; cross-region posture (Aurora Global Database,
global tables) justified by a written RPO/RTO.
\chk DynamoDB: access patterns enumerated \emph{before} keys were chosen; on-demand vs.\
provisioned decided from measured traffic; GSIs sparse where possible; TTL understood as
asynchronous.
\chk Redshift: grain declared for every fact table; distribution and sort keys validated with
\texttt{EXPLAIN} and \texttt{SVV\_TABLE\_INFO} against production-shaped data volumes.
\chk WLM configured with queue monitoring rules; concurrency scaling enabled for bursty BI.
\chk COPY loads use manifests and multiple input files; MERGE-based upserts are idempotent under
re-run.
\chk Serverless capacity bounds (max RPU, max ACU) set so a runaway query cannot run away with
the budget.
\chk Materialized views have an owner and a refresh cadence; stale-MV staleness is visible to
consumers.
\end{itemize}

\section{ETL and Orchestration (Glue, EMR, MWAA, Step Functions; Chapters 9--12)}
\index{checklist!ETL}%
\begin{itemize}
\chk Every job re-runnable without duplicates: bookmarks, watermarks, or idempotent MERGE writes.
\chk EMR clusters are transient for recurring batch; Spot only on task nodes; durable state lives
in S3 via EMRFS, never on cluster HDFS.
\chk Glue jobs parameterized (no hardcoded dates or bucket names); Flex execution for non-urgent
jobs.
\chk DAGs deployed with \texttt{catchup=False} unless a backfill is intended; every alert task has
\texttt{trigger\_rule="one\_failed"} (or the state-machine equivalent).
\chk Sensors tuned: \texttt{poke\_interval} sized to the expected arrival cadence.
\chk Large data moves through S3, never through XComs or state-machine payloads.
\chk Failure path proven, not assumed: a forced failure fires the alert and pages the right queue.
\chk Schema contracts enforced at ingestion (registry compatibility checks), not at query time.
\chk CDC lag monitored on \texttt{CDCLatencySource}/\texttt{CDCLatencyTarget}, not on task status.
\end{itemize}

\section{Streaming (Kinesis, Firehose, MSK, Flink, Lambda; Chapters 13--15)}
\index{checklist!streaming}%
\begin{itemize}
\chk Shard/partition count derived from both throughput limits (bytes \emph{and} records per
second) plus headroom; hot-key risk assessed in partition-key choice.
\chk \texttt{IteratorAgeMilliseconds} (or consumer-lag equivalent) alarmed well before retention
expiry.
\chk Consumers idempotent: dedup on a unique event ID or upsert by key; poison records route to a
DLQ that someone watches.
\chk Firehose streams have source-record backup enabled and an alarm on
\texttt{DeliveryToRedshift.Success} (or destination equivalent); the intermediate S3 prefix has
lifecycle rules.
\chk Lambda transforms are stateless; anything needing windows, dedup, or joins runs in Flink,
not in a delivery transform.
\chk Flink jobs: every stateful operator bounded (window end, join interval, TTL); savepoint taken
before any redeploy; RocksDB state size trended.
\chk Watermark policy stated in one sentence (``no event older than X'') and agreed with
downstream consumers.
\chk Replay drill performed: reprocess from a checkpoint/savepoint without double-landing output.
\end{itemize}

\section{Machine Learning and MLOps (Chapters 16--20)}
\index{checklist!ML}%
\begin{itemize}
\chk Training datasets built through governed access (Lake Formation grants), so PII masks apply
to extraction; no broad S3 policies on ML roles.
\chk Training query reviewed column-by-column for target leakage (``is this knowable before the
outcome?'').
\chk Objective metric matches the business cost (e.g., F1 for imbalanced churn), not accuracy by
default.
\chk One feature-group definition feeds both offline training and online serving---the skew
defense.
\chk Deployment gated on Model Registry approval; every production model is versioned, attributable,
and rollback-ready.
\chk Model Monitor baselines refreshed on a cadence tied to retraining; drift alerts triaged by a
human before any automatic retrain.
\chk Managed Spot training used with checkpointing to S3 so interruption resumes, not restarts.
\chk AI-API pipelines (Textract/Comprehend) idempotent on object identity so retries never
re-bill; raw OCR/PII output lives in a restricted zone with field-level redaction downstream.
\end{itemize}

\section{Security, Governance, and Compliance (Chapters 8, 21--23)}
\index{checklist!security}%
\begin{itemize}
\chk Root user: MFA on, no access keys, no daily use; all operations through Identity Center or
roles.
\chk \texttt{IAMAllowedPrincipals} revoked on every governed database and table; Lake Formation
grants verified from a consumer role, not the admin role.
\chk Permissions boundaries on delegated role creation; SCPs cap what any account---including its
admins---may do.
\chk Private connectivity: gateway endpoints for S3/DynamoDB, interface endpoints with endpoint
policies elsewhere; no public route for data-plane traffic.
\chk KMS key policies restrict decryption to organization principals, closing the exfiltration
path that survives network controls.
\chk CloudTrail data events enabled for lake buckets and delivered to an independent log-archive
account.
\chk Secrets in Secrets Manager with automatic rotation; zero long-lived keys in CI (OIDC
federation instead).
\chk Data-quality gates halt the pipeline on hard failures and report on thresholds
(\texttt{Completeness >= 0.95 with threshold 0.9} pattern); quality metrics trended in CloudWatch.
\chk Lineage consultable per dataset: given a bad batch, you can name its consumers and its
origin change in minutes.
\chk Macie findings feed an automated quarantine (Object-Locked bucket), with an audit record
linking finding ID to action; custom identifiers tuned for precision before rollout.
\chk Erasure designed, not improvised: crypto-shredding keys per subject where re-identification
is not required, or an enumerated-deletion runbook that includes replicas, snapshots, and streams.
\end{itemize}

\section{Cost and Operations (Chapters 0, 23, 24)}
\index{checklist!cost}%
\begin{itemize}
\chk Budgets with 50/80/100\% and forecast alerts on every account that can spend money.
\chk Every resource tagged (\texttt{Environment}, \texttt{Project}, \texttt{Owner}); untagged
spend reviewed weekly.
\chk Per-TB-scanned workloads (Athena, Spectrum) enforced through partitioning + Parquet;
bytes-scanned reported per query in reviews.
\chk Idle-cost audit automated: running clusters, attached-but-stopped databases, NAT gateways,
old snapshots, and high base-capacity Serverless endpoints.
\chk Terraform state remote (S3 + DynamoDB locking); \texttt{main} always deployable; PR checks
prove \texttt{plan} + \texttt{dbt build} before merge; Prod deploys require named approval.
\chk Schema changes follow expand-contract with a compatibility shim and a deferred drop.
\end{itemize}

\begin{tipbox}
A checklist item is done when you can point to the evidence---a config output, a metric graph, a
drill log. If the honest answer is ``we believe so,'' the box stays empty.
\end{tipbox}
